%! Absolute Value Equations \& Inequalities — Unit 1, Lesson 1.2 — Lesson Plan
% Gradual-release plan (user direction, 2026-08-31; extended course-wide 2026-09-01): warm-up →
% guided notes (with guided practice) → individual practice → debrief → close. There is no group
% activity; the "you do" is done alone. No tiers, no exit ticket; homework is in-packet.
\documentclass[10pt]{article}
\usepackage{algebra2-article}
\usepackage{algebra2-boxes}
\usepackage{graphicx}   % -article does NOT load graphicx; needed for thumbnails/screenshots
\graphicspath{{images/}}
\newcommand{\TallMath}[1]{$\displaystyle #1\rule[-1.4em]{0pt}{3.2em}$}
\newcommand{\CourseName}{Algebra 2}
\newcommand{\MeetingLength}{60 minutes}
\newcommand{\UnitNumberName}{Unit 1: Linear Functions \quad}
\newcommand{\LessonNumberName}{Lesson 1.2: Absolute Value Equations \& Inequalities}

\begin{document}
\begin{center}
    {\Large\bfseries \CourseName \\
    {\small\bfseries \UnitNumberName \LessonNumberName}}
\end{center}

\begin{tcolorbox}[colback=forestbg, colframe=forest, boxrule=0.9pt, arc=2mm,
    left=2mm, right=2mm, top=1.5mm, bottom=1.5mm]
    \textbf{Primary Objective:} Students read $|x-h|$ as the \emph{distance} between $x$ and $h$,
    and use that meaning to \emph{create}, \emph{solve}, and \emph{interpret} absolute value
    equations and inequalities in one variable. They \emph{isolate} the bars before splitting into
    two cases, handle the $c=0$ and $c<0$ special cases from the distance meaning rather than from a
    memorised rule, distinguish $|X|<c$ (one interval, \emph{between}) from $|X|>c$ (two rays,
    \emph{outside}), and represent a solution set in set notation, in interval notation, and on a
    number line. Lesson~1.1 wrote lines that cross zero exactly \emph{once}; today's rule reaches
    the same value \emph{twice}.
    \\[2pt]
    \textbf{Standards (2023 VA SOL):} This lesson delivers the whole of \textbf{A2.EI.1} ---
    represent, solve, and interpret absolute value equations and inequalities in one variable:
    \textbf{A2.EI.1a} (create an absolute value \emph{equation} to model a context),
    \textbf{A2.EI.1b} (solve an absolute value equation algebraically and verify),
    \textbf{A2.EI.1c} (create an absolute value \emph{inequality} to model a context),
    \textbf{A2.EI.1d} (solve an absolute value inequality and represent the solution set with set
    notation, interval notation, and a number line), and \textbf{A2.EI.1e} (verify solutions, and
    explain and interpret solutions in context).
    \\[2pt]
    \textbf{Lesson model:} \emph{Gradual release} --- I do, we do, you do alone. A warm-up
    activates prior knowledge; \textbf{Guided Notes} build the distance meaning with the class on
    one worked context (a radio tower on a parkway) and close with a \emph{Guided Practice} box (the
    ``we do''); an \textbf{Individual Practice} block on the same page is the ``you do'' --- three
    problems, worked silently and alone; a \textbf{debrief} consolidates and surfaces the errors;
    \textbf{homework} extends the practice. \emph{There is no group activity in this lesson} ---
    today's independent work is done individually.
\end{tcolorbox}

\renewcommand{\arraystretch}{1.4}
\begin{skillbox}[Priority Ideas \& Skills]{goldbox}
\begin{tabularx}{\linewidth}{@{}X|X@{}}
\textbf{Priority Ideas \& Skills} & \textbf{Key Understandings (the ``why'')} \\[2pt]
\hline
\vspace{-6pt}
\begin{itemize}[leftmargin=1.1em, nosep, after=\vspace{2pt}]
  \item Read $|x-h|$ as the \textbf{distance} between $x$ and $h$.
  \item \textbf{Isolate} the bars, then apply the two-case rule.
  \item Solve $|X|=c$ for $c>0$, $c=0$, and $c<0$.
  \item Solve $|X|<c$ (\emph{and}) and $|X|>c$ (\emph{or}).
  \item Write a set three ways: set, interval, number line.
  \item Model a ``within so much of'' context and interpret the answer.
\end{itemize}
&
\vspace{-6pt}
\begin{itemize}[leftmargin=1.1em, nosep, after=\vspace{2pt}]
  \item A distance is never negative, and \emph{two} points sit the same distance from a
        centre --- one on each side. That single fact generates every rule today.
  \item You cannot split until the bars stand alone, because the two-case rule is a statement
        about $|X|$, not about the equation around it.
  \item \textbf{Target misconception:} that $\le$ and $\ge$ produce ``nearly the same'' answer.
        They produce \emph{opposite} sets --- one interval versus two rays.
  \item Set, interval, and number-line notation are three spellings of one set; the endpoint
        symbol ($[$ vs.\ $($, filled vs.\ open) is the only thing carrying inclusion.
\end{itemize}
\end{tabularx}
\end{skillbox}

\begin{skillbox}[{Vocabulary, Concepts, \& Theorems}]{sky}
\renewcommand{\arraystretch}{1.3}
\begin{tabularx}{\linewidth}{@{}l X@{}}
\textbf{Absolute value as distance} & $|x-h|$ is how far $x$ is from $h$ on the number line; never negative. \\
\textbf{Two-case rule} & $|X|=c$ with $c>0$ splits into $X=c$ or $X=-c$. $c=0$ gives one solution; $c<0$ gives none. \\
\textbf{Compound ``and''} & $|X|<c$ becomes $-c<X<c$ --- one interval, \emph{between} the endpoints. \\
\textbf{Compound ``or''} & $|X|>c$ becomes $X<-c$ or $X>c$ --- two rays, \emph{outside} the endpoints. \\
\textbf{Set notation} & $\{x \mid 7 \le x \le 17\}$ --- ``all $x$ such that\dots'' \\
\textbf{Interval notation} & $[7,17]$; brackets include an endpoint, parentheses exclude it. Next to $\pm\infty$, always a parenthesis. \\
\textbf{No solution / all reals} & $|X|=c$ or $|X|<c$ with $c<0$ has none; $|X|\ge c$ with $c<0$ holds for every $x$. \\
\end{tabularx}
\end{skillbox}

\begin{fixedskillbox}[Lesson at a Glance --- \MeetingLength]{forestbg}
\renewcommand{\arraystretch}{1.25}
\begin{tabularx}{\linewidth}{@{}l c X X@{}}
\textbf{Phase} & \textbf{Min} & \textbf{Students} & \textbf{Teacher} \\
\hline
Warm-Up & 5 & Three prior skills, alone then check with a neighbour & Circulate; debrief item~3 aloud (``could a shaded picture ever come in two pieces?'') and leave it hanging \\
Guided Notes & 20 & Fill the vocabulary and the four sections with the class, then \emph{Guided Practice} together & \emph{I do / we do}: run the whole block on \emph{one} context, the parkway; three \texttt{work} blocks at the board \\
Individual Practice & 15 & \emph{Alone and quiet}: three problems --- an equation, an ``and'' inequality, an ``or'' inequality & Circulate; questions, cues, prompts --- \emph{not} answers. Note who to call on in the debrief \\
Debrief & 10 & Correct their own pages in a second colour & Work all three individual-practice problems on the board; land ``one interval or two rays'' \\
Close \& Homework & 10 & Start the homework page in class; preview 1.3 & Launch item~1 aloud, then circulate; name what changed today \\
\end{tabularx}
\end{fixedskillbox}

\begin{fixedskillbox}[Warm-Up --- Activate Prior Knowledge \& Spiral Review]{forestbg}
The Warm-Up rehearses exactly what today leans on, without naming anything. \textbf{(1)} Evaluating
$|-7|$, $|4|$, $|-3|+|5|$, $-|-6|$ re-establishes that the bars measure size, and the follow-up
(\emph{how many numbers have absolute value 7?}) is the seed of the whole \emph{two-case rule}.
\textbf{(2)} Reading distances off a number line and writing them with bars --- $|11-3|$,
$|6-(-2)|$, and finally the open $|x-6|$ --- is the exact move notes section~1 formalises.
\textbf{(3)} Shading $x\ge-1$ and ``between $-3$ and $2$'' rehearses filled vs.\ open dots, and the
closing question --- \emph{could a shaded picture ever come in two separate pieces?} --- is
deliberately left hanging until notes section~3 answers it. \textbf{Debrief item 3 aloud, briefly},
and land only this: \emph{both of your pictures were one piece --- hold on to the question of
whether that always happens.} \textbf{This page is set at 12pt} --- larger type than the rest of the
packet, deliberately.
\end{fixedskillbox}

\begin{skillbox}[Hook]{forestbg}
    ``A straight parkway is marked in miles, $0$ to $20$. A ranger's radio tower stands at
    \textbf{mile 12}, and a handheld radio reaches it only from \textbf{within 5 miles}.'' Point at
    the picture and ask: \textbf{``Two hikers are standing exactly 5 miles from the tower. Where
    are they?''} (Miles $7$ and $17$.) Then the question that carries the lesson:
    \textbf{``Why are there two of them, and not one?''} (One on each side.) Do not name anything
    yet --- get the picture and the two answers on the board, then say: \emph{every question today
    is this picture. How far is $x$ from the tower --- and is that close enough?}
\end{skillbox}

\begin{skillbox}[Guided Notes --- the Lesson (20 min)]{forestbg}
\begin{multicols}{2}
\textbf{1. Distance --- so it answers twice.} Name $|x-h|$ as the distance between $x$ and $h$; on
the parkway a hiker at mile $x$ is $|x-12|$ from the tower. Work the first \texttt{work} block:
$|x-12|=5$ splits to $x=17$ or $x=7$ --- the hook's own two markers, falling out of the algebra.
Verify both by substituting back; that substitution \emph{is} what A2.EI.1b means by verifying.

\textbf{2. Isolate before you split.} Fill the three-row table ($c>0$, $c=0$, $c<0$) reading each
row off the distance meaning, not a rule. Then the \texttt{work} block on $3|x+1|-2=10$, and the
sentence that matters: \emph{splitting at the first line gives $3x+1-2=10$ --- one answer, and a
wrong one.}

\columnbreak
\textbf{3. ``Less th\emph{AND}'' vs.\ ``Great\emph{OR}.''} The two-column table, side by side, on
the \emph{same} $|x-12|$ and the \emph{same} $5$: $\le$ shades $7\le x\le17$ as one piece; $\ge$
shades two. Then the gold caution box --- test the hiker at \textbf{mile 3}, distance $9$: is
$9\le5$? No. Is $9\ge5$? Yes. \emph{The wrong symbol does not shift the answer a little; it hands
you the opposite set.} This is the target misconception, and this box is where it dies.

\textbf{4. Three ways to write one set.} $7\le x\le17$ as $\{x \mid \dots\}$, as $[7,17]$, and
shaded. Brackets/filled include, parentheses/open exclude, and $\infty$ always takes a parenthesis
because it is not a number you can reach.

\textbf{Guided Practice (the ``we do'').} $2|x-5|-3=7$ --- isolate, split ($x=10$ or $x=0$), then
change one symbol twice and watch the answer change shape. Do this one \emph{with} the class ---
it is the last thing they see before working alone.
\end{multicols}
\end{skillbox}

\begin{skillbox}[Individual Practice --- \emph{On Your Own} (15 min, silent)]{redbox}
\textbf{Launch (1 min):} ``Same three moves you just watched --- now nobody helps you. Three
problems. When you get stuck, look back at the notes above, not at your neighbour.'' The block sits
on the last page of the Guided Notes, so nothing new is handed out.
\begin{multicols}{2}
\textbf{What students do.} \textbf{Problem~1} is the equation with something in the way:
$4|x+1|-5=15$, isolated to $|x+1|=5$, split to $x=4$ or $x=-6$, then \emph{verified}.
\textbf{Problem~2} is the ``and'' inequality, $|2x-1|\le5 \Rightarrow -2\le x\le3$, written in all
three notations including the shaded line. \textbf{Problem~3} is the deliberate partner,
$|3x-6|\ge9 \Rightarrow x\le-1$ or $x\ge5$ --- same shape of expression, opposite shape of answer,
and the first place they must write a union and two $\infty$ parentheses. Together the three cover
the equation, the ``and'', and the ``or''.

\columnbreak
\textbf{What the teacher does.} Circulate the whole 15 minutes; keep it silent. Give a
\emph{question, cue, or prompt}, never an answer:
\begin{itemize}[leftmargin=1.1em, nosep]
  \item 1: ``What is stuck to the bars? Undo that first --- what is $|x+1|$ by itself?''
  \item 1: ``You have one answer. The bars measure a distance --- where is the other side?''
  \item 2: ``Is this \emph{closer than} or \emph{farther than}? So between, or outside?''
  \item 3: ``Does your picture have a gap in the middle? Then it needs two intervals and a $\cup$.''
  \item 3: ``What goes next to $\infty$ --- a bracket or a parenthesis? Why?''
  \item Finished early: ``Change problem~3's $\ge$ to $\le$. Which numbers switch sides?''
\end{itemize}
Note names as you go: one student for each of the three problems, to put work on the board in the
debrief.
\end{multicols}
\end{skillbox}

\boxguard
\begin{skillbox}[Debrief (10 min --- what goes on the board, in a second color)]{forestbg}
Students correct their own Individual Practice in a second colour as you go. Work all three, but
protect the time for items~2 and~3.
\begin{enumerate}[leftmargin=1.3em, nosep, itemsep=2pt]
  \item \textbf{Problem~1 --- isolate, then split.} Put the four lines up in order:
        $4|x+1|=20$, $|x+1|=5$, $x+1=5$ \emph{or} $x+1=-5$, $x=4$ or $x=-6$. Ask who split at line
        one, and show what they got instead. Then verify $x=4$: $4|5|-5=15$.
  \item \textbf{The must-land moment: one interval or two rays.} Put problems~2 and~3 side by side
        on the board with their number lines. Say it out loud --- \emph{closer than gives one
        connected piece; farther than gives two} --- and point back at the mile-3 hiker from the
        notes' caution box. This is the parkway's $|x-12|\le5$ versus $|x-12|\ge5$ all over again.
  \item \textbf{Problem~2 --- dividing a compound inequality.} $-5\le2x-1\le5$ becomes
        $-4\le2x\le6$ becomes $-2\le x\le3$: whatever you do, you do to \emph{all three} parts.
        Check by testing $x=0$: $|{-1}|\le5$, true.
  \item \textbf{Problem~3 --- the union and the two parentheses.}
        $(-\infty,-1]\cup[5,\infty)$. Ask why the brackets face the way they do, and why $\infty$
        never gets one. Test $x=0$ (not a solution) and $x=6$ ($|12|\ge9$, true).
\end{enumerate}
If time is short, cut item~1's verification and protect items~2 and~4.
\end{skillbox}

\boxguard
\begin{skillbox}[Active Monitoring --- Watch For]{redbox}
\begin{itemize}[leftmargin=1.2em, nosep]
  \item \textbf{splitting before the bars are alone} --- $3|x+1|-2=10$ becoming $3x+1-2=10$
        (notes~2, Individual Practice~1, homework~1b) --- the single most common error today;
  \item giving only the positive case and stopping at one answer (Warm-Up~1, notes~1);
  \item \textbf{treating $\le$ and $\ge$ as nearly the same answer} (Individual Practice~2 vs.~3,
        homework~2) --- the lesson's target misconception. Send them to the mile-3 hiker, not to
        the rule;
  \item writing an ``or'' answer as a single chain, $5\le x\le-1$, which reads as empty ---
        make them shade it and see the gap (Individual Practice~3);
  \item solving $|x-4|=-3$ by dropping the bars to get $x=1$, instead of reading it as an
        impossible distance (homework~4a);
  \item at the boundary $k=0$ in homework~4c: ``distance $<0$'' is still impossible, so $k=0$
        belongs to the no-solution set;
  \item bracketing $\infty$, or dividing only part of a compound inequality (homework~2, 3b).
\end{itemize}
Cold-call prompts: ``How far is that from the centre?'' ``Closer than, or farther than?'' ``One
piece or two --- convince me from the picture.'' ``Is that endpoint in or out? Which symbol says so?''
\end{skillbox}

\boxguard
\begin{skillbox}[Reinforcement \& Extension]{goldbox}
\textbf{Homework} (in the packet, collected next class): two equations, the second needing isolation
first; a deliberate contrast pair ($|x+3|<4$ against $|x+3|>4$) with the ``why one piece, why two''
question; \emph{reading an inequality off a shaded number line} --- the third representation, and
the only item that runs backwards; the special cases plus the $k\le0$ boundary; a thermostat
tolerance model with an interpret-the-answer follow-up; and an SOL-style multiple-choice item on
$|x+2|>3$. \textbf{Extension:} building $|x-10|\le6$ from the endpoints $4$ and $16$ (find the
centre first), and why miles $7$ and $17$ belong to \emph{both} the in-range and out-of-range
answers. \textbf{DeltaMath override:} well covered there --- swap in a DeltaMath set on absolute
value equations and inequalities if you prefer auto-graded practice. \textbf{Preview:}
Lesson~1.3, \emph{Absolute Value Functions \& Transformations} --- today you \emph{solved}
$|ax+b|=c$; next you \emph{graph} $y=|ax+b|$, and the two solutions you found become the two places
a V crosses a horizontal line.
\end{skillbox}

% ── Teacher notes — one per component, in packet order ──
\begin{teachernote}[Warm-Up]
Item~1's follow-up --- \emph{how many different numbers have an absolute value of 7?} --- is the
whole lesson in one question; if the room says ``one,'' you have found your hook and section~1 will
land hard. Item~2's last blank, $|x-6|$, is the first time they write a distance with a
\emph{variable} in it, and it is exactly the expression notes section~1 needs; do not rush past it.
Item~3's closing question is deliberately left unresolved --- \emph{could a shaded picture ever come
in two separate pieces?} Let them answer ``no'' or shrug, and do \emph{not} correct it. Notes
section~3 pays it off, and the surprise is worth more than the five seconds you would save.
\textbf{This page is set at 12pt} --- larger type than the rest of the packet, deliberately, so the
five minutes go on the thinking rather than on decoding it. It is still one page; if you add to it,
check that it stays one page.
\end{teachernote}

\begin{teachernote}[Guided Notes]
Twenty minutes, paced by section: 3 for the vocabulary and hook, 3 for section~1, 4 for section~2,
5 for section~3, 3 for section~4, and 2 for the Guided Practice box worked \emph{with} the class.
\textbf{Keep the whole notes block on one context} --- the parkway --- so ``distance from the
tower'' is doing the work in every section rather than four unrelated procedures. Section~3 is the
one to protect: the two-column table and its gold caution box are where the lesson's target
misconception dies, and the mile-3 hiker ($|3-12|=9$, so $9\le5$ is false but $9\ge5$ is true) is
the concrete case that makes ``opposite sets'' undeniable. Six \texttt{work} blocks are printed on
these pages (the $|x-12|=5$ split, the $3|x+1|-2=10$ isolation, the Guided Practice box, and three
more in Individual Practice); students write in the reserved space while the key prints the same
steps there, so the two files paginate alike. If you fall behind, compress section~4's endpoint
table to one sentence --- but do \emph{not} cut section~2's ``splitting at the first line'' warning;
Individual Practice~1 and homework~1b both target it.
\end{teachernote}

\begin{teachernote}[Individual Practice]
Fifteen minutes, and \textbf{keep it silent} --- this is the only block all period where a student
finds out what they can do without help. Budget roughly 5 minutes each. The three problems are
deliberately the three \emph{shapes} of answer: an equation (two points), an ``and'' (one interval),
and an ``or'' (two rays). A student who can do 1 and 2 but not 3 has not accepted that an answer can
come in two pieces --- send them back to their own Warm-Up~3 picture and to the notes' caution box,
not to a rule. Watch for the ``or'' answer written as a single chain, $5\le x\le-1$: ask them to
shade it and the contradiction shows itself. Fast finishers change problem~3's $\ge$ to $\le$ and
report which numbers switch sides; that is the work you want on the board in the debrief. Common
slips: splitting $4|x+1|-5=15$ before isolating; dividing only part of $-5\le2x-1\le5$; and putting
a bracket next to $\infty$.
\end{teachernote}

\begin{teachernote}[Homework]
Item~1(b) is the isolation check --- expect $5x+2-1=14$ from anyone who did not internalise
section~2. Item~3 is the only item that runs \emph{backwards}, from a picture to an inequality, and
it is the best single predictor of whether the distance meaning stuck: they must find the centre
($3$) and the reach ($4$) before they can write $|x-3|\le4$. Item~4(c) is the trap --- $k=0$ is
\emph{included} in the no-solution answer, because ``distance $<0$'' is still impossible; students
who write $k<0$ have the idea but not the boundary. Item~6 is the formative check: sort into ``chose
B (got it) / chose A (solved the $<$ version --- the target misconception, still alive) / chose C
(kept one ray, forgot the other side) / chose D (confused $>$ with $\ge-3$)'' to decide whether
Lesson~1.3 opens with a one-minute re-look at ``less th\emph{AND}'' versus ``great\emph{OR}.'' Three
\texttt{work} blocks (items 1b, 5b, and one in the extension). With the group activity cut, the last
10 minutes are for \emph{starting} this page in class --- launch item~1 aloud, then circulate; the
rest is due next class.
\end{teachernote}

\end{document}
